\scanpage{74}
\mspage{67}
% archive sheet number: 73
Prouvons la partie c) du th.-scholie, partie $1^\circ$. Cela résulte de a), une fois vu que pour $s\in\unsure{I|\mathbb{I}}$, $\deleted{\unsure{s\backslash}}I'_{/s}$ satisfait aux $=$s hypothèses que $I'$. Or $I$ \deleted{\unsure{étant|est}} ordonné, $\deleted{\unsure{s\backslash}}I_{/s}$ est une sous-catégorie ouverte de $I$, donc $I'_{/s}$ une sous-cat ouverte de $I'$, donc \inserted{$(I'_{/s})^\circ$} un sous-ens ordonné fermé de $I'^\circ$. Ce dernier étant noethérien et sobre, il en est de $=$ des \unsure{dits} fermés, d'où la conclusion.

Prouvons la partie $2^\circ)$ du th.-scholie. Considérons
\diagram{
\begin{tikzcd}[ampersand replacement=\&,column sep=large,row sep=large]
M(I')_c\arrow[r,hook]\&M(I')\&M(I')_f\arrow[l,hook]\arrow[d,"u_*^f"']\\
M(I)_c\arrow[r,hook]\arrow[u,"u_c^*"']\&M(I)\arrow[u,"u^*"']\&M(I)_f\arrow[l,hook]
\end{tikzcd}
}
,
où \deleted{$\unsure{u^*}$} \inserted{$u^*$ et a fortiori $u_c^*$} est compatible avec $W(I),W(I')$, et par hyp. $u_*$ est \unsure{comp.|comp} aux $W(I')_f$, $W(I)_f$. En passant aux cat. de fractions, on trouve
